% Introduction (Task 42): three RQs and the contributions of handoff/37_narrative_v2.md section 4.
\section{Introduction}
\label{sec:intro}

Machine-learning network intrusion detection systems (NIDS) classify network flows from statistical features, and transformer models for tabular data have
been applied to this task \cite{gorishniy2021,manocchio2024}.
Whether attention is a faithful explanation is debated in natural language processing
\cite{jain2019,wiegreffe2019,serrano2019}, and the question has received less attention for flow-feature models.

Two further questions arise in a security setting.
First, an attacker who perturbs a flow to evade the detector may also change the explanation, so the behaviour of
attention under attack matters for any analyst who reads it \cite{ghorbani2019,dombrowski2019,pruthi2020}.
Second, if attention and attribution methods agree on clean traffic but disagree on adversarial traffic, their
consistency could serve as an attack signal \cite{fidel2020}.
A detector built on such a signal has to be evaluated against attackers that model it \cite{tramer2020,carlini2017},
and the strength of that attacker decides what the evaluation can show.
Both questions only make sense for attacks whose flows are consistent with written protocol constraints, so we
restricted the attacker to a hand-built model of flow-feature constraints and checked every adversarial row with a
separate validity checker \cite{pierazzi2020,alhussien2024}.

Our central finding concerns the second question.
Against attackers that optimised the exact score of the detector they faced, with surrogate gradients only for the
\nRankScores{} rank-based consistency scores, adding attention-attribution consistency scores to a
confidence-and-distance detector raised the share of successful and undetected attack rows, by \axHeadVoneSeventeenAbs{}
on 2017 and by \axHeadVoneEighteenAbs{} on 2018, with \axHeadOppSeedsPhrase{} of the opposite sign.
An attacker with differentiable stand-ins for the same scores had shown a reduction of \axStandinVoneSeventeenAbs{} on
2017, so that result did not hold.
Giving the attacker of the base detector the same iterate selection as the stronger attacker of the extended detector
changed no row, so the reversal does not come from the selection rule.

We asked three research questions (RQ).
\begin{description}
  \item[RQ1] Is attention a faithful explanation of the FT-Transformer, on clean flows and under protocol-constrained
             attacks?
  \item[RQ2] Do attention-attribution consistency scores improve a confidence-and-distance detector, first against an
             attacker that ignores them and then against attackers that model them, from differentiable stand-ins to
             the detector's exact score? Is any gain specific to attention?
  \item[RQ3] How much do the answers depend on the dataset and its data quality?
\end{description}
We answered them with an FT-Transformer trained with \nSeeds{} seeds on two datasets, LYCOS-IDS2017 and
LycoS-Unicas-IDS2018, using the \nCommonClasses{} attack classes common to both for every comparison.
The contributions are the following.
\begin{enumerate}
  \item \emph{Faithfulness of attention in NIDS.} On clean data, a comparison of three attention variants with four
        attribution methods by agreement and by deletion and insertion faithfulness on two datasets; under
        protocol-constrained attacks, measurements showing that attention was stiff but less faithful than Integrated
        Gradients (IG) and occlusion, and that the faithfulness loss after an evasion follows the prediction flip
        (RQ1; Section~\ref{sec:rq1}).
  \item \emph{An adaptive evaluation of consistency monitoring with three attacker levels.} An attacker that ignores
        the consistency scores (A6), one with differentiable stand-ins (A8) and one that optimises the exact detector
        score with surrogate gradients only for the rank-based scores (A8x), compared under an equal-effort protocol
        with outcome rules recorded before the runs.
        The gain of the extended detector against the stand-in attacker reversed against the exact-objective attacker:
        on 2017 the difference went from \axStandinVoneSeventeen{} to \axHeadVoneSeventeen.
        This is a concrete case in which a stand-in adaptive attack overstated what a detector achieves
        (RQ2; Sections~\ref{sec:protocol} and~\ref{sec:rq2}).
  \item \emph{Control detectors that separate attention from other scores.} \nControlsAll{} controls, \nControlsAttacked{} of them
        attacked by adversaries that modelled their own scores.
        The gain against the ignoring attacker was largely obtained without attention, and no control was clearly
        better than the confidence-and-distance detector on both datasets under its own adaptive attacker
        (RQ2; Section~\ref{sec:rq2}).
\end{enumerate}
The constraint model with its validity checker and the near-duplicate analysis are method details that support these
results (Sections~\ref{sec:constraints} and~\ref{sec:rq3}).
Both halves of the title, stable but not faithful attention and the limits of consistency monitoring, refer to these two
datasets only.
